% =====================================================================
%  CHƯƠNG 8. HỆ ĐA TÁC TỬ HỖ TRỢ MUA SẮM (Tutorial MAS, PyTorch trên GPU)
% =====================================================================
\chapter{Hệ đa tác tử hỗ trợ mua sắm}\label{chap:8}

Chương này hiện thực hệ mua sắm của Tutorial MAS \cite{ch8:tutorial} (tác tử ý định là mạng PyTorch trên GPU) để đo
luận điểm của Chương~1: phân vai giúp xử lý lỗi, nhưng mỗi chặng thông điệp đều có giá về độ trễ.

\section{Đặc tả PEAS và luồng thông điệp}
Catalog tổng hợp có \SL{ch8_shop/data/n_products} mặt hàng sáu loại; kho có \SL{ch8_shop/data/in_stock} mặt hàng còn
hàng. Môi trường (E) là phiên mua sắm; mọi tác tử chỉ \emph{đọc} dữ liệu của mình (Bảng~\ref{tab:ch8-peas}).

\begin{table}[!htbp]\centering
\caption{Đặc tả PEAS của năm thành phần trong hệ mua sắm.}\label{tab:ch8-peas}
\small
\begin{tabularx}{\linewidth}{L{2.9cm}XL{3.6cm}}
\toprule
\textbf{Tác tử} & \textbf{Cảm biến (S) $\to$ Hành động (A), thông điệp} & \textbf{Hiệu năng (P)} \\
\midrule
IntentAgent & câu người dùng $\to$ CLASSIFY/INTENT \{ý định, độ tin cậy\}; dưới ngưỡng $\tau$ thì hỏi lại
  & accuracy, macro-F1 \\
CatalogAgent & câu, catalog $\to$ SEARCH/CANDIDATES \{loại, hãng, ngân sách, id\}; không nhận ra loại thì trả rỗng
  & đúng tập ứng viên \\
Recommendation & ứng viên $\to$ RANK/RANKED: lọc giá sau khuyến mãi, xếp theo đánh giá, top-3 & không vượt ngân sách \\
InventoryAgent & id, kho $\to$ CHECK\_STOCK/STOCK \{số lượng, tuổi dữ liệu\} & dữ liệu mới ($\le 60$ s) \\
Coordinator & mọi phản hồi $\to$ gọi theo thứ tự, kiểm hợp đồng, thử lại & đúng, không khẳng định sai, trễ \\
\bottomrule
\end{tabularx}
\end{table}

\hinhtikz{ch8_flow}{Luồng thông điệp của hệ mua sắm: Coordinator gửi bốn yêu cầu (nét liền) và nhận bốn phản hồi
(nét đứt) theo thứ tự 1--4; lời gọi kho lỗi được thử lại tối đa hai lần.}{ch8-flow}

Mọi lời gọi đi qua \texttt{Coordinator.call} (Đoạn mã~\ref{code:shop_mas__Coordinator.call}). Phản hồi chỉ được nhận
khi đúng lớp \texttt{Message}, đúng người gửi/nhận, đúng loại, đủ trường; riêng STOCK còn phải có tuổi dữ liệu
$\le 60$~s. Lỗi được ghi rồi thử lại; hết lượt thì trả \texttt{chua\_xac\_minh}, \emph{không} khẳng định còn hàng.
Độ trễ theo đồng hồ ảo: mỗi chặng trễ lognormal trung vị 2~ms cộng thời gian tính thật, quá hạn tốn trọn 80~ms.

\noindent\begin{minipage}{\linewidth}
\maDai{shop_mas__Coordinator.call}{Coordinator gửi yêu cầu, kiểm hợp đồng phản hồi và thử lại khi lỗi.}
\end{minipage}

\section{Bộ phân loại ý định bằng PyTorch trên GPU}
Câu truy vấn được sinh tất định từ 14 khuôn câu Việt--Anh mỗi ý định, thêm biến thể như người dùng thật (hậu tố hội
thoại, bỏ dấu 30\,\%, hoán đổi ký tự). Mấu chốt là \textbf{chia khuôn câu trước khi sinh câu} (9 train, 2 val, 3
test), nên test gồm những \emph{cách nói chưa từng gặp}: \SL{ch8_shop/data/train} câu train,
\SL{ch8_shop/data/val} val, \SL{ch8_shop/data/test} test, \SL{ch8_shop/data/n_intents} ý định.

Mỗi câu là một túi đặc trưng (từ, cặp từ, 3-gram ký tự, ở cả dạng có dấu và không dấu) băm CRC32 vào $2^{16}$ ô theo
tinh thần fastText \cite{ch8:fasttext}; \texttt{IntentNet} (Đoạn mã~\ref{code:shop_intent__IntentNet}) lấy trung bình
embedding rồi qua hai lớp tuyến tính, tổng \SL{ch8_shop/train/params} tham số. Huấn luyện AdamW \cite{ch8:pytorch} trên
\textbf{\SL{ch8_shop/train/device}}, \SL{ch8_shop/train/epochs} epoch trong \SL{ch8_shop/train/train_seconds}~s,
giữ trọng số epoch \SL{ch8_shop/train/best_epoch} có val accuracy cao nhất.

\begin{figure}[H]\centering
\includegraphics[width=\linewidth]{ch8_learning.pdf}
\caption{Đường học của IntentNet: train gần 100\,\% sau vài epoch, val dừng thấp hơn rõ rệt -- khoảng cách đo đúng
độ khó của khuôn câu mới.}\label{fig:ch8-learning}
\end{figure}

\maDai{shop_intent__IntentNet}{Mạng phân loại ý định: EmbeddingBag trên đặc trưng băm và đầu phân loại hai lớp.}



\begin{figure}[H]\centering
\begin{minipage}[c]{0.50\linewidth}\centering
  {\small\setlength{\tabcolsep}{5pt}\begin{adjustbox}{max width=\linewidth}\begin{tabular}{lrrr}
\toprule
Ý định & P & R & F1 \\
\midrule
Tìm SP & 0{,}761 & 0{,}922 & 0{,}834 \\
So sánh & 0{,}732 & 1{,}000 & 0{,}845 \\
Tồn kho & 0{,}819 & 0{,}656 & 0{,}728 \\
Giá/KM & 0{,}957 & 0{,}989 & 0{,}973 \\
Chào hỏi & 0{,}818 & 0{,}684 & 0{,}745 \\
Khác & 1{,}000 & 0{,}714 & 0{,}833 \\
\midrule
\textbf{Acc / Macro-F1} &  & 0{,}832 & 0{,}826 \\
\bottomrule
\end{tabular}
\end{adjustbox}}\par\smallskip
  {\small (a) Precision (P), recall (R), F1}
\end{minipage}\hfill
\begin{minipage}[c]{0.48\linewidth}\centering
  \includegraphics[width=\linewidth]{ch8_confusion.pdf}\par
  {\small (b) Ma trận nhầm lẫn}
\end{minipage}
\caption{Bộ phân loại ý định trên tập test gồm các khuôn câu chưa gặp khi huấn luyện.}\label{fig:ch8-report}
\end{figure}

\nhanxet{Accuracy test \SL{ch8_shop/test/accuracyPct}\,\% (macro-F1 \SL{ch8_shop/test/macro_f1}), thấp hơn hẳn train vì
test đổi cách nói. Yếu nhất là \emph{hỏi tồn kho} (recall \SL{ch8_shop/test/per_class/hoi_ton_kho/recallPct}\,\%, hay
nhầm sang \emph{so sánh}) và \emph{chào hỏi} (recall \SL{ch8_shop/test/per_class/chao_hoi/recallPct}\,\%),
Hình~\ref{fig:ch8-report}b. Vì vậy cần ngưỡng tin cậy, chọn trên val: $\tau = \SL{ch8_shop/train/tau}$. Hệ nhận
\SL{ch8_shop/test/coveragePct}\,\% câu test, đúng \SL{ch8_shop/test/accepted_accPct}\,\% trong số đó; còn lại được hỏi
lại thay vì đoán sai.}

\section{Demo, kiểm thử và mô phỏng lỗi}
Hộp dưới là nhật ký thông điệp thật của một yêu cầu demo; lần gọi kho đầu được cài sẵn lỗi quá hạn nên Coordinator
chờ hết 80~ms, ghi sự kiện, gửi lại và nhận tồn kho hợp lệ. \SL{ch8_shop/tests/run} kiểm thử đơn vị (phân tích ngân
sách, câu lạ, ngân sách, hàng hết, ba kiểu lỗi kho, sai người gửi, pipeline tin dữ liệu cũ) đều đạt.

\begin{ketqua}{Nhật ký thông điệp của Coordinator (\texttt{scripts/ch8\_run.py})}
{\footnotesize\textbf{>>> Tìm giúp mình laptop dưới 20 triệu}\par
[~~~2.5~ms]~coordinator~->~intent~~~~~~CLASSIFY~~~~~\{'query':~'Tìm~giúp~mình~laptop~dư\par
[~~11.1~ms]~~~~~~intent~->~coordinator~INTENT~~~~~~~\{'intent':~'tim\_san\_pham',~'conf':\par
[~~14.2~ms]~coordinator~->~catalog~~~~~SEARCH~~~~~~~\{'query':~'Tìm~giúp~mình~laptop~dư\par
[~~17.1~ms]~~~~~catalog~->~coordinator~CANDIDATES~~~\{'category':~'laptop',~'brands':~[\par
[~~18.9~ms]~coordinator~->~recommend~~~RANK~~~~~~~~~\{'ids':~['p01',~'p02',~'p03',~'p04\par
[~~22.0~ms]~~~recommend~->~coordinator~RANKED~~~~~~~\{'ids':~['p04',~'p07',~'p05']\}\par
[~~23.5~ms]~coordinator~->~inventory~~~CHECK\_STOCK~~\{'ids':~['p04',~'p07',~'p05']\}\par
[~103.5~ms]~~!!~inventory:~inventory~không~trả~lời~trước~hạn~chót~(lần~1)\par
[~105.3~ms]~coordinator~->~inventory~~~CHECK\_STOCK~~\{'ids':~['p04',~'p07',~'p05']\}\par
[~106.9~ms]~~~inventory~->~coordinator~STOCK~~~~~~~~\{'stock':~\{'p04':~8,~'p07':~1,~'p0\par
\textit{=> ok p04, p07, p05 (9 thông điệp, 106{,}9 ms)}
}
\end{ketqua}

Để cô lập phần điều phối, cả hai hệ nhận \emph{ý định đúng} của \SL{ch8_shop/data/test} câu test và phản hồi được so
với đáp án chuẩn. Với xác suất 0, 10 hoặc 30\,\%, kho gặp một trong ba lỗi đều nhau: quá hạn, thiếu trường
\texttt{stock}, hoặc bản chụp cũ 3.600~s lệch kho thật ở khoảng một phần ba mặt hàng. Pipeline đơn khối cùng logic
nhưng gọi hàm trực tiếp, không kiểm hợp đồng, không thử lại.

\hinhpdf{ch8_faults}{1}{Mô phỏng lỗi kho trên \SL{ch8_shop/data/test} yêu cầu: MAS có Coordinator so với pipeline đơn
khối; nhãn trắng là số lần khẳng định sai tồn kho, trục trễ theo thang log.}{ch8-faults}

\nhanxet{Khi kho không lỗi hai hệ cùng đúng \SL{ch8_shop/faults/r0/mas/completionPct}\,\%, nhưng MAS tốn
\SL{ch8_shop/faults/r0/mas/msgs_mean} thông điệp, \SL{ch8_shop/faults/r0/mas/lat_mean}~ms mỗi yêu cầu so với
\SL{ch8_shop/faults/r0/pipe/lat_mean}~ms: giá của phân vai. Ở mức lỗi 30\,\%, pipeline còn
\SL{ch8_shop/faults/r30/pipe/completionPct}\,\%, \textbf{\SL{ch8_shop/faults/r30/pipe/unsafe} lần khẳng định sai} tồn
kho và \SL{ch8_shop/faults/r30/pipe/errors} yêu cầu hỏng; MAS giữ \SL{ch8_shop/faults/r30/mas/completionPct}\,\%,
không khẳng định sai lần nào, \SL{ch8_shop/faults/r30/mas/unverified} yêu cầu báo ``chưa xác minh''. Có bộ phân
loại ở đầu chuỗi, tỉ lệ đúng toàn trình là \SL{ch8_shop/e2e/completionPct}\,\%: nút thắt là hiểu câu.}
